% PN Junction Physics Template
% Topics: Built-in potential, depletion region, I-V characteristics, capacitance, breakdown
% Style: Semiconductor device physics with computational analysis

\documentclass[a4paper, 11pt]{article}
\usepackage[utf8]{inputenc}
\usepackage[T1]{fontenc}
\usepackage{amsmath, amssymb}
\usepackage{graphicx}
\usepackage{siunitx}
\usepackage{booktabs}
\usepackage{subcaption}
\usepackage[makestderr]{pythontex}
\usepackage{geometry}
\geometry{margin=1in}
\usepackage{hyperref}

% Theorem environments
\newtheorem{definition}{Definition}[section]
\newtheorem{theorem}{Theorem}[section]
\newtheorem{example}{Example}[section]
\newtheorem{remark}{Remark}[section]

\title{PN Junction Physics: Electrostatic Analysis and Current-Voltage Characteristics}
\author{Semiconductor Device Laboratory}
\date{\today}

\begin{document}
\maketitle

\begin{abstract}
This report presents a comprehensive computational analysis of PN junction physics,
covering the formation of the depletion region, built-in potential, electrostatic
field profiles, current-voltage characteristics governed by the Shockley diode
equation, junction capacitance (both depletion and diffusion), and breakdown mechanisms.
We analyze silicon PN junctions with varying doping concentrations, compute key parameters
including built-in voltage ($V_{bi} = 0.7$--$0.9$ V), saturation current ($I_s \sim 10^{-15}$ A),
and breakdown voltage ($V_{BR} \sim 10$--$100$ V), and examine the transition from thermal
equilibrium to forward/reverse bias conditions.
\end{abstract}

\section{Introduction}

The PN junction forms the foundation of modern semiconductor devices, from diodes and
bipolar junction transistors to solar cells and light-emitting diodes. When P-type
(acceptor-doped) and N-type (donor-doped) semiconductors are brought into contact,
carrier diffusion establishes a depletion region with an internal electric field that
balances diffusion and drift currents at equilibrium.

\begin{definition}[PN Junction]
A PN junction is a metallurgical junction between P-type semiconductor (excess holes,
$N_A$ acceptors/cm$^3$) and N-type semiconductor (excess electrons, $N_D$ donors/cm$^3$),
characterized by a built-in potential $V_{bi}$ and depletion width $W$ at thermal equilibrium.
\end{definition}

\section{Theoretical Framework}

\subsection{Built-in Potential}

At thermal equilibrium, the Fermi level is constant across the junction. The built-in
potential barrier arises from the difference in quasi-Fermi levels before contact.

\begin{theorem}[Built-in Voltage]
For a silicon PN junction at temperature $T$ with intrinsic carrier concentration $n_i$:
\begin{equation}
V_{bi} = \frac{kT}{q} \ln\left(\frac{N_A N_D}{n_i^2}\right)
\end{equation}
where $k = 1.38 \times 10^{-23}$ J/K is Boltzmann's constant, $q = 1.6 \times 10^{-19}$ C
is the elementary charge, and $n_i = 1.5 \times 10^{10}$ cm$^{-3}$ for silicon at 300 K.
\end{theorem}

\subsection{Depletion Region Analysis}

\begin{definition}[Depletion Width]
Under reverse bias $V_R$ or zero bias, the depletion region extends into both P and N sides:
\begin{equation}
W = \sqrt{\frac{2\epsilon_s}{q}\left(\frac{1}{N_A} + \frac{1}{N_D}\right)(V_{bi} + V_R)}
\end{equation}
where $\epsilon_s = 11.7 \epsilon_0 = 1.04 \times 10^{-12}$ F/cm is the silicon permittivity.
\end{definition}

The electric field profile is triangular, peaking at the metallurgical junction:
\begin{equation}
E_{max} = \frac{qN_A x_p}{\epsilon_s} = \frac{qN_D x_n}{\epsilon_s}
\end{equation}

\subsection{Current-Voltage Characteristics}

\begin{theorem}[Shockley Diode Equation]
The current through an ideal PN junction under applied voltage $V$ is:
\begin{equation}
I = I_s\left[\exp\left(\frac{qV}{nkT}\right) - 1\right]
\end{equation}
where $I_s$ is the saturation current (reverse-bias leakage) and $n$ is the ideality
factor ($n = 1$ for ideal diode, $n = 1.5$--$2$ for real diodes with recombination).
\end{theorem}

The saturation current depends on minority carrier properties:
\begin{equation}
I_s = qA\left(\frac{D_n n_{p0}}{L_n} + \frac{D_p p_{n0}}{L_p}\right)
\end{equation}
where $D_{n,p}$ are diffusion coefficients, $L_{n,p}$ are diffusion lengths, and
$n_{p0}, p_{n0}$ are equilibrium minority carrier concentrations.

\subsection{Junction Capacitance}

\begin{definition}[Depletion Capacitance]
The voltage-dependent depletion capacitance per unit area:
\begin{equation}
C_j = \frac{\epsilon_s}{W} = \sqrt{\frac{q\epsilon_s}{2(V_{bi} + V_R)}\left(\frac{N_A N_D}{N_A + N_D}\right)}
\end{equation}
\end{definition}

\begin{definition}[Diffusion Capacitance]
Under forward bias, minority carrier storage contributes:
\begin{equation}
C_d = \frac{qI}{nkT}\tau_t
\end{equation}
where $\tau_t$ is the transit time (minority carrier lifetime).
\end{definition}

\subsection{Breakdown Mechanisms}

\begin{remark}[Breakdown Voltage]
At high reverse bias, two mechanisms cause breakdown:
\begin{itemize}
\item \textbf{Avalanche breakdown}: Impact ionization when $E_{max} \approx 10^6$ V/cm (larger $V_{BR}$ for lightly doped junctions)
\item \textbf{Zener breakdown}: Band-to-band tunneling for heavily doped junctions ($V_{BR} < 5$ V)
\end{itemize}
\end{remark}

\section{Computational Analysis}

\begin{pycode}
import numpy as np
import matplotlib.pyplot as plt
from scipy.constants import k as k_B, e as q_e, epsilon_0
from scipy.optimize import fsolve

np.random.seed(42)

# Physical constants
k = k_B  # Boltzmann constant (J/K)
q = q_e  # Elementary charge (C)
T = 300  # Temperature (K)
Vt = k * T / q  # Thermal voltage (V)
eps_0 = epsilon_0  # Vacuum permittivity (F/m)
eps_r_Si = 11.7  # Relative permittivity of silicon
eps_s = eps_r_Si * eps_0  # Silicon permittivity (F/m)
n_i = 1.5e10  # Intrinsic carrier concentration (cm^-3) at 300K

# Convert permittivity to F/cm for semiconductor units
eps_s_cgs = eps_s * 1e-2  # F/cm

# Doping concentrations (cm^-3)
N_A = 1e16  # Acceptor concentration (P-side)
N_D = 1e15  # Donor concentration (N-side)

# Built-in potential
V_bi = Vt * np.log((N_A * N_D) / (n_i**2))

# Depletion widths at equilibrium
def depletion_width(V_applied):
    """Calculate total depletion width"""
    V_total = V_bi - V_applied  # Subtract for forward, add for reverse
    if V_total < 0:
        return 0  # No depletion in strong forward bias (not physical for this model)
    W = np.sqrt((2 * eps_s_cgs / q) * ((1/N_A) + (1/N_D)) * V_total)
    return W * 1e4  # Convert cm to micrometers

def depletion_widths_separate(V_applied):
    """Calculate P-side and N-side depletion widths"""
    W_total = depletion_width(V_applied) * 1e-4  # Back to cm
    x_p = W_total * N_D / (N_A + N_D)
    x_n = W_total * N_A / (N_A + N_D)
    return x_p * 1e4, x_n * 1e4  # Convert to micrometers

# Electric field profile
def electric_field_profile(V_applied):
    """Generate electric field profile across junction"""
    x_p, x_n = depletion_widths_separate(V_applied)
    x_p_cm = x_p * 1e-4
    x_n_cm = x_n * 1e-4

    E_max = (q * N_A * x_p_cm) / eps_s_cgs  # V/cm

    # Position arrays
    x_p_array = np.linspace(-x_p, 0, 50)
    x_n_array = np.linspace(0, x_n, 50)

    # Field in P-side (linear from 0 to E_max)
    E_p = E_max * (1 + x_p_array / x_p)

    # Field in N-side (linear from E_max to 0)
    E_n = E_max * (1 - x_n_array / x_n)

    return x_p_array, x_n_array, E_p, E_n, E_max

# Current-voltage characteristics
def diode_current(V, I_s, n=1):
    """Shockley diode equation"""
    return I_s * (np.exp(V / (n * Vt)) - 1)

# Estimate saturation current (simplified)
# Assume D_n = 25 cm^2/s, D_p = 10 cm^2/s, L_n = L_p = 10 um
D_n = 25  # cm^2/s
D_p = 10  # cm^2/s
L_n = 10e-4  # cm
L_p = 10e-4  # cm
n_p0 = n_i**2 / N_A
p_n0 = n_i**2 / N_D
A = 1e-4  # Junction area (cm^2)

I_s = q * A * (D_n * n_p0 / L_n + D_p * p_n0 / L_p)

# Junction capacitance
def junction_capacitance(V_applied):
    """Depletion capacitance per unit area (F/cm^2)"""
    V_total = V_bi - V_applied
    if V_total <= 0:
        return np.inf
    C_j = np.sqrt((q * eps_s_cgs / (2 * V_total)) * ((N_A * N_D) / (N_A + N_D)))
    return C_j

def diffusion_capacitance(V_forward, tau_t=1e-9):
    """Diffusion capacitance (F)"""
    I = diode_current(V_forward, I_s, n=1)
    if I <= 0:
        return 0
    C_d = (q * I * tau_t) / (k * T)
    return C_d

# Voltage arrays
V_reverse = np.linspace(-5, 0, 100)
V_forward = np.linspace(0, 0.8, 100)
V_full = np.linspace(-5, 0.8, 200)

# Depletion width vs voltage
W_vs_V = [depletion_width(V) for V in V_reverse]

# I-V curve
I_reverse = [diode_current(V, I_s, n=1) for V in V_reverse]
I_forward = [diode_current(V, I_s, n=1) for V in V_forward]
I_full = [diode_current(V, I_s, n=1) for V in V_full]

# Capacitance vs voltage
C_j_vs_V = [junction_capacitance(V) for V in V_reverse]
C_d_vs_V = [diffusion_capacitance(V) for V in V_forward if V > 0]

# Create comprehensive figure
fig = plt.figure(figsize=(16, 12))

# Plot 1: Built-in potential and band diagram (conceptual)
ax1 = fig.add_subplot(3, 3, 1)
x_junction = np.linspace(-2, 2, 200)
E_c = np.where(x_junction < 0, 0.56, 0.56 - V_bi)
E_v = np.where(x_junction < 0, -0.56, -0.56 - V_bi)
E_f = np.zeros_like(x_junction)
ax1.plot(x_junction, E_c, 'b-', linewidth=2, label='$E_c$')
ax1.plot(x_junction, E_v, 'r-', linewidth=2, label='$E_v$')
ax1.plot(x_junction, E_f, 'k--', linewidth=1.5, label='$E_F$')
ax1.axvline(x=0, color='gray', linestyle=':', alpha=0.7)
ax1.set_xlabel('Position ($\\mu$m)')
ax1.set_ylabel('Energy (eV)')
ax1.set_title(f'Energy Band Diagram ($V_{{bi}}$ = {V_bi:.3f} V)')
ax1.legend(fontsize=9)
ax1.grid(True, alpha=0.3)
ax1.set_ylim(-1, 1)

# Plot 2: Electric field profile at equilibrium
ax2 = fig.add_subplot(3, 3, 2)
x_p_eq, x_n_eq, E_p_eq, E_n_eq, E_max_eq = electric_field_profile(0)
ax2.plot(x_p_eq, E_p_eq, 'b-', linewidth=2.5, label='P-side')
ax2.plot(x_n_eq, E_n_eq, 'r-', linewidth=2.5, label='N-side')
ax2.axhline(y=0, color='black', linewidth=0.5)
ax2.axvline(x=0, color='gray', linestyle='--', alpha=0.7)
ax2.fill_between(x_p_eq, 0, E_p_eq, alpha=0.2, color='blue')
ax2.fill_between(x_n_eq, 0, E_n_eq, alpha=0.2, color='red')
ax2.set_xlabel('Position ($\\mu$m)')
ax2.set_ylabel('Electric Field (V/cm)')
ax2.set_title(f'Electric Field Profile (Equilibrium, $E_{{max}}$ = {E_max_eq:.2e} V/cm)')
ax2.legend(fontsize=9)
ax2.grid(True, alpha=0.3)

# Plot 3: Depletion width vs reverse bias
ax3 = fig.add_subplot(3, 3, 3)
ax3.plot(V_reverse, W_vs_V, 'b-', linewidth=2.5)
ax3.set_xlabel('Reverse Bias Voltage (V)')
ax3.set_ylabel('Depletion Width ($\\mu$m)')
ax3.set_title('Depletion Width vs Reverse Bias')
ax3.grid(True, alpha=0.3)
ax3.set_xlim(-5, 0)

# Plot 4: I-V characteristics (linear scale)
ax4 = fig.add_subplot(3, 3, 4)
I_full_mA = np.array(I_full) * 1e3  # Convert to mA
ax4.plot(V_full, I_full_mA, 'b-', linewidth=2.5)
ax4.axhline(y=0, color='black', linewidth=0.5)
ax4.axvline(x=0, color='gray', linestyle='--', alpha=0.7)
ax4.set_xlabel('Voltage (V)')
ax4.set_ylabel('Current (mA)')
ax4.set_title(f'Diode I-V Characteristics ($I_s$ = {I_s:.2e} A)')
ax4.grid(True, alpha=0.3)
ax4.set_xlim(-5, 0.8)

# Plot 5: I-V characteristics (semi-log scale)
ax5 = fig.add_subplot(3, 3, 5)
I_forward_pos = np.array([max(i, 1e-18) for i in I_forward])
V_forward_array = np.array(V_forward)
ax5.semilogy(V_forward_array, I_forward_pos * 1e3, 'r-', linewidth=2.5, label='$n=1$')
I_forward_n2 = [diode_current(V, I_s, n=1.5) for V in V_forward]
I_forward_n2_pos = np.array([max(i, 1e-18) for i in I_forward_n2])
ax5.semilogy(V_forward_array, I_forward_n2_pos * 1e3, 'b--', linewidth=2, label='$n=1.5$')
ax5.set_xlabel('Forward Voltage (V)')
ax5.set_ylabel('Current (mA, log scale)')
ax5.set_title('Forward Characteristics (Semi-log Plot)')
ax5.legend(fontsize=9)
ax5.grid(True, alpha=0.3, which='both')
ax5.set_xlim(0, 0.8)

# Plot 6: Junction capacitance vs reverse bias
ax6 = fig.add_subplot(3, 3, 6)
C_j_vs_V_pF = np.array(C_j_vs_V) * A * 1e12  # Convert to pF
ax6.plot(V_reverse, C_j_vs_V_pF, 'b-', linewidth=2.5)
ax6.set_xlabel('Reverse Bias Voltage (V)')
ax6.set_ylabel('Junction Capacitance (pF)')
ax6.set_title('Depletion Capacitance vs Reverse Bias')
ax6.grid(True, alpha=0.3)
ax6.set_xlim(-5, 0)

# Plot 7: 1/C^2 vs V (for V_bi extraction)
ax7 = fig.add_subplot(3, 3, 7)
inv_C_sq = 1 / (np.array(C_j_vs_V_pF)**2)
ax7.plot(V_reverse, inv_C_sq, 'b-', linewidth=2.5)
# Linear fit to extract V_bi
from scipy.stats import linregress
slope, intercept, r_val, p_val, std_err = linregress(V_reverse, inv_C_sq)
V_bi_extracted = -intercept / slope
ax7.plot(V_reverse, slope * V_reverse + intercept, 'r--', linewidth=2,
         label=f'Fit: $V_{{bi}}$ = {V_bi_extracted:.3f} V')
ax7.set_xlabel('Voltage (V)')
ax7.set_ylabel('1/$C_j^2$ (pF$^{-2}$)')
ax7.set_title('1/$C^2$ vs V (Built-in Voltage Extraction)')
ax7.legend(fontsize=9)
ax7.grid(True, alpha=0.3)

# Plot 8: Electric field for different biases
ax8 = fig.add_subplot(3, 3, 8)
bias_values = [0, -1, -2, -3]
colors_bias = plt.cm.viridis(np.linspace(0, 0.9, len(bias_values)))
for i, V_bias in enumerate(bias_values):
    x_p_b, x_n_b, E_p_b, E_n_b, E_max_b = electric_field_profile(V_bias)
    ax8.plot(x_p_b, E_p_b, color=colors_bias[i], linewidth=2)
    ax8.plot(x_n_b, E_n_b, color=colors_bias[i], linewidth=2,
             label=f'$V$ = {V_bias} V')
ax8.axhline(y=0, color='black', linewidth=0.5)
ax8.axvline(x=0, color='gray', linestyle='--', alpha=0.7)
ax8.set_xlabel('Position ($\\mu$m)')
ax8.set_ylabel('Electric Field (V/cm)')
ax8.set_title('Electric Field vs Bias Voltage')
ax8.legend(fontsize=8)
ax8.grid(True, alpha=0.3)

# Plot 9: Doping concentration profile
ax9 = fig.add_subplot(3, 3, 9)
x_doping = np.linspace(-3, 3, 500)
# Abrupt junction model
N_profile = np.where(x_doping < 0, N_A, -N_D)
ax9.semilogy(x_doping, np.abs(N_profile), 'b-', linewidth=2.5)
ax9.axvline(x=0, color='gray', linestyle='--', alpha=0.7)
ax9.fill_between(x_doping[x_doping < 0], 1e10, N_A, alpha=0.3, color='red', label='P-type')
ax9.fill_between(x_doping[x_doping >= 0], 1e10, N_D, alpha=0.3, color='blue', label='N-type')
ax9.set_xlabel('Position ($\\mu$m)')
ax9.set_ylabel('Doping Concentration (cm$^{-3}$)')
ax9.set_title(f'Doping Profile ($N_A$ = {N_A:.1e}, $N_D$ = {N_D:.1e})')
ax9.legend(fontsize=9)
ax9.grid(True, alpha=0.3, which='both')
ax9.set_ylim(1e13, 1e18)

plt.tight_layout()
plt.savefig('pn_junction_analysis.pdf', dpi=150, bbox_inches='tight')
plt.close()

# Store computed values for table
W_0 = depletion_width(0)
x_p_0, x_n_0 = depletion_widths_separate(0)
E_max_0 = electric_field_profile(0)[4]
C_j_0 = junction_capacitance(0) * A * 1e12  # pF
I_0p7 = diode_current(0.7, I_s, n=1) * 1e3  # mA

\end{pycode}

\begin{figure}[htbp]
\centering
\includegraphics[width=\textwidth]{pn_junction_analysis.pdf}
\caption{Comprehensive PN junction analysis for silicon with $N_A = 10^{16}$ cm$^{-3}$
and $N_D = 10^{15}$ cm$^{-3}$: (a) Energy band diagram at thermal equilibrium showing
built-in potential barrier of \py{f"{V_bi:.3f}"} V; (b) Triangular electric field profile
with peak field \py{f"{E_max_0:.2e}"} V/cm at the metallurgical junction; (c) Square-root
dependence of depletion width on reverse bias voltage; (d) Current-voltage characteristics
on linear scale showing exponential forward conduction and saturation reverse current
\py{f"{I_s:.2e}"} A; (e) Semi-logarithmic I-V plot comparing ideal ($n=1$) and non-ideal
($n=1.5$) diode behavior; (f) Voltage-dependent junction capacitance decreasing with
reverse bias; (g) Linear relationship in $1/C^2$ vs $V$ plot enabling extraction of
built-in voltage; (h) Electric field expansion under increasing reverse bias; (i) Abrupt
doping concentration profile with P-side acceptor density and N-side donor density.}
\label{fig:pn_junction}
\end{figure}

\section{Results}

\subsection{Junction Parameters at Thermal Equilibrium}

\begin{pycode}
print(r"\begin{table}[htbp]")
print(r"\centering")
print(r"\caption{Computed PN Junction Parameters (T = 300 K)}")
print(r"\begin{tabular}{lcc}")
print(r"\toprule")
print(r"Parameter & Value & Units \\")
print(r"\midrule")
print(f"Acceptor concentration ($N_A$) & {N_A:.1e} & cm$^{{-3}}$ \\\\")
print(f"Donor concentration ($N_D$) & {N_D:.1e} & cm$^{{-3}}$ \\\\")
print(f"Built-in voltage ($V_{{bi}}$) & {V_bi:.4f} & V \\\\")
print(f"Thermal voltage ($V_t$) & {Vt:.4f} & V \\\\")
print(f"Depletion width (equilibrium) & {W_0:.4f} & $\\mu$m \\\\")
print(f"P-side depletion ($x_p$) & {x_p_0:.4f} & $\\mu$m \\\\")
print(f"N-side depletion ($x_n$) & {x_n_0:.4f} & $\\mu$m \\\\")
print(f"Maximum electric field ($E_{{max}}$) & {E_max_0:.2e} & V/cm \\\\")
print(f"Junction capacitance ($C_j$) & {C_j_0:.2f} & pF \\\\")
print(f"Saturation current ($I_s$) & {I_s:.2e} & A \\\\")
print(r"\bottomrule")
print(r"\end{tabular}")
print(r"\label{tab:parameters}")
print(r"\end{table}")
\end{pycode}

\subsection{Current at Different Voltages}

\begin{pycode}
print(r"\begin{table}[htbp]")
print(r"\centering")
print(r"\caption{Diode Current vs Applied Voltage}")
print(r"\begin{tabular}{ccc}")
print(r"\toprule")
print(r"Voltage (V) & Current ($n=1$) & Current ($n=1.5$) \\")
print(r"\midrule")

voltage_points = [-5, -1, 0, 0.5, 0.6, 0.7, 0.8]
for V in voltage_points:
    I_n1 = diode_current(V, I_s, n=1)
    I_n15 = diode_current(V, I_s, n=1.5)
    if abs(I_n1) < 1e-6:
        print(f"{V:.1f} & {I_n1:.2e} A & {I_n15:.2e} A \\\\")
    else:
        print(f"{V:.1f} & {I_n1*1e3:.2f} mA & {I_n15*1e3:.2f} mA \\\\")

print(r"\bottomrule")
print(r"\end{tabular}")
print(r"\label{tab:current}")
print(r"\end{table}")
\end{pycode}

\subsection{Depletion Width vs Reverse Bias}

\begin{pycode}
print(r"\begin{table}[htbp]")
print(r"\centering")
print(r"\caption{Depletion Width Under Reverse Bias}")
print(r"\begin{tabular}{cccc}")
print(r"\toprule")
print(r"$V_R$ (V) & $W$ ($\mu$m) & $x_p$ ($\mu$m) & $x_n$ ($\mu$m) \\")
print(r"\midrule")

reverse_voltages = [0, -1, -2, -3, -5]
for V_R in reverse_voltages:
    W = depletion_width(V_R)
    x_p, x_n = depletion_widths_separate(V_R)
    print(f"{V_R:.0f} & {W:.4f} & {x_p:.4f} & {x_n:.4f} \\\\")

print(r"\bottomrule")
print(r"\end{tabular}")
print(r"\label{tab:depletion}")
print(r"\end{table}")
\end{pycode}

\section{Discussion}

\begin{example}[Asymmetric Depletion]
For this junction with $N_A/N_D = 10$, the depletion region extends approximately
10 times farther into the lightly-doped N-side than the P-side. This is evident from:
\[
\frac{x_n}{x_p} = \frac{N_A}{N_D} = 10
\]
The asymmetry is critical for breakdown voltage design in power diodes.
\end{example}

\begin{remark}[Forward Voltage Drop]
The forward voltage required to achieve significant current ($I \gg I_s$) is approximately:
\begin{pycode}
V_forward_10mA = Vt * np.log(10e-3 / I_s)
print(f"At $I = 10$ mA, $V_f \\approx {V_forward_10mA:.3f}$ V.")
\end{pycode}
This $\sim 0.7$ V threshold is characteristic of silicon diodes and is widely used in
circuit design.
\end{remark}

\begin{example}[Breakdown Voltage Estimation]
The critical field for avalanche breakdown in silicon is approximately $E_{crit} \approx 3 \times 10^5$ V/cm.
Using the maximum field relation:
\begin{pycode}
E_crit = 3e5  # V/cm
# E_max = sqrt(2*q*N*(V_bi + V_R)/eps_s)
# For one-sided junction (N_D << N_A), use N_D
V_BR_estimated = (E_crit**2 * eps_s_cgs) / (2 * q * N_D) - V_bi
print(f"Estimated breakdown voltage: $V_{{BR}} \\approx {V_BR_estimated:.1f}$ V.")
\end{pycode}
Higher breakdown voltages require lower doping concentrations.
\end{example}

\begin{remark}[Temperature Dependence]
The built-in voltage decreases with temperature at approximately $-2$ mV/K due to
the $T$ dependence in the logarithm and the exponential increase of $n_i^2 \propto T^3 \exp(-E_g/kT)$.
\end{remark}

\section{Conclusions}

This computational analysis demonstrates the fundamental physics of PN junctions:

\begin{enumerate}
\item The built-in potential for this silicon junction is $V_{bi} = \py{f"{V_bi:.3f}"}$ V,
determined by the doping ratio and intrinsic carrier concentration.

\item The depletion width at equilibrium is $W_0 = \py{f"{W_0:.4f}"}$ $\mu$m, expanding
to $\py{f"{depletion_width(-5):.4f}"}$ $\mu$m under $-5$ V reverse bias.

\item The saturation current is $I_s = \py{f"{I_s:.2e}"}$ A, setting the leakage
current floor in reverse bias.

\item Forward bias reduces the barrier, enabling exponential current growth with voltage
according to the Shockley equation.

\item The junction capacitance decreases from $\py{f"{C_j_0:.2f}"}$ pF at zero bias,
following $C_j \propto (V_{bi} + V_R)^{-1/2}$ under reverse bias.

\item The estimated breakdown voltage for this doping level is approximately
$\py{f"{V_BR_estimated:.1f}"}$ V, suitable for low-voltage rectification applications.
\end{enumerate}

The computational framework enables rapid exploration of junction design parameters,
essential for optimizing diodes, solar cells, and bipolar transistors.

\section*{References}

\begin{itemize}
\item Sze, S. M. \& Ng, K. K. \textit{Physics of Semiconductor Devices}, 3rd ed. Wiley, 2006.
\item Streetman, B. G. \& Banerjee, S. K. \textit{Solid State Electronic Devices}, 7th ed. Pearson, 2014.
\item Neamen, D. A. \textit{Semiconductor Physics and Devices: Basic Principles}, 4th ed. McGraw-Hill, 2011.
\item Schroder, D. K. \textit{Semiconductor Material and Device Characterization}, 3rd ed. Wiley-IEEE Press, 2005.
\item Pierret, R. F. \textit{Semiconductor Device Fundamentals}. Addison-Wesley, 1996.
\item Muller, R. S. \& Kamins, T. I. \textit{Device Electronics for Integrated Circuits}, 3rd ed. Wiley, 2002.
\item Tyagi, M. S. \textit{Introduction to Semiconductor Materials and Devices}. Wiley, 1991.
\item Brennan, K. F. \& Brown, A. S. \textit{Theory of Modern Electronic Semiconductor Devices}. Wiley, 2002.
\item Tsividis, Y. \textit{Operation and Modeling of the MOS Transistor}, 3rd ed. Oxford University Press, 2010.
\item Ashcroft, N. W. \& Mermin, N. D. \textit{Solid State Physics}. Brooks/Cole, 1976.
\item Kittel, C. \textit{Introduction to Solid State Physics}, 8th ed. Wiley, 2004.
\item Shur, M. \textit{Physics of Semiconductor Devices}. Prentice Hall, 1990.
\item Rhoderick, E. H. \& Williams, R. H. \textit{Metal-Semiconductor Contacts}, 2nd ed. Oxford University Press, 1988.
\item Anderson, B. L. \& Anderson, R. L. \textit{Fundamentals of Semiconductor Devices}, 2nd ed. McGraw-Hill, 2004.
\item Hu, C. \textit{Modern Semiconductor Devices for Integrated Circuits}. Prentice Hall, 2009.
\item Taur, Y. \& Ning, T. H. \textit{Fundamentals of Modern VLSI Devices}, 2nd ed. Cambridge University Press, 2009.
\item Green, M. A. \textit{Solar Cells: Operating Principles, Technology and System Applications}. Prentice Hall, 1982.
\item Roulston, D. J. \textit{An Introduction to the Physics of Semiconductor Devices}. Oxford University Press, 1999.
\item Singh, J. \textit{Semiconductor Devices: Basic Principles}. Wiley, 2000.
\item Fonash, S. J. \textit{Solar Cell Device Physics}, 2nd ed. Academic Press, 2010.
\end{itemize}

\end{document}
